\clearpage
\renewcommand{\sectioncolor}{bio}
\section{The question, and the premise checked}
\markboth{The premise}{}

\begin{keybox}[title={The question as posed}]
\textit{Since our mathematics came mostly from human physical realities of thought, our mathematics
is limited to our cognitive limitations. In that case, how do we invent new mathematics to explore
the new worlds of nano/quantum-scale biological systems inside cells, molecules, atoms? What can we
do about that?}
\end{keybox}

\noindent It is the natural inference to draw from Part VI --- and it is the one Lakoff \& Núñez
themselves are careful not to draw. Separating the part that is correct from the part that does not
follow is what tells you what to do.

\subsection{Two claims that get conflated}

\begin{center}
\begin{tikzpicture}[font=\footnotesize,
  c/.style={rounded corners=4pt,draw=#1,line width=1.2pt,fill=#1!7,inner sep=9pt,
            text width=7.35cm,align=left,minimum height=3.4cm}]
 \node[c=bio] (t) at (-4.3,0)
   {\textbf{\color{bio}\normalsize CONCEIVABILITY CONSTRAINT}\\[4pt]
    You must be able to \emph{build} the concept out of grounded pieces, via metaphor and blending.\\[5pt]
    \textcolor{bio}{\textbf{TRUE.}} This is the book's actual thesis.\\[3pt]
    {\scriptsize It constrains the \textbf{route}.}};
 \node[c=crimson] (f) at (4.3,0)
   {\textbf{\color{crimson}\normalsize IMAGINABILITY CONSTRAINT}\\[4pt]
    You must be able to \emph{picture} the result.\\[5pt]
    \textcolor{crimson}{\textbf{FALSE}} --- and nothing in Part VI asserts it.\\[3pt]
    {\scriptsize It would constrain the \textbf{destination}. It does not exist.}};
 \node[font=\Large\bfseries,text=slate] at (0,0) {$\neq$};
\end{tikzpicture}
\end{center}

\noindent The evidence against the second is overwhelming, and you already hold it. Nobody pictures
an infinite-dimensional separable Hilbert space. Nobody pictures a non-commutative torus. Nobody
pictures an $\infty$-groupoid. Mathematicians compute with all three fluently and without error.

\begin{ideabox}[title={How the escape actually works: substituting the source domain}]
When vision and folk physics run out, mathematics does not become disembodied. It \textbf{re-grounds
in a different body of experience}:
\begin{center}
\begin{tikzpicture}[font=\footnotesize,
  g/.style={rounded corners=3pt,draw=#1,line width=0.9pt,fill=#1!8,inner sep=6pt,
            text width=3.55cm,align=center,minimum height=1.45cm}]
 \node[g=slate]   (a) at (0,0)     {\textbf{seeing a shape}\\{\scriptsize vision, folk geometry}};
 \node[g=slate]   (b) at (4.15,0)  {\textbf{moving an object}\\{\scriptsize motor, collection, container}};
 \node[g=nano]    (c) at (8.30,0)  {\textbf{manipulating marks}\\{\scriptsize symbol pushing, rewriting}};
 \node[g=quantum] (d) at (12.45,0) {\textbf{chasing a diagram}\\{\scriptsize path-following, pattern match}};
 \draw[slate!45,line width=0.8pt,dashed] (2.1,-1.0)--(2.1,1.0);
 \draw[slate!45,line width=0.8pt,dashed] (6.25,-1.0)--(6.25,1.0);
 \draw[slate!45,line width=0.8pt,dashed] (10.4,-1.0)--(10.4,1.0);
 \node[font=\scriptsize,text=crimson,align=center] at (6.25,-1.42)
   {\textbf{visualisability ends here} --- but the grounding does not};
\end{tikzpicture}
\end{center}
\vspace{2pt}
Hand, eye and motor cortex are still doing the work in the right-hand two columns. They are simply
not doing \emph{spatial imagery}. This distinction is not a technicality: it tells you exactly which
lever to pull. If the obstacle were imagination, you would be stuck. Because the obstacle is
grounding, you can \textbf{build a new grounding}.
\end{ideabox}

\noindent And the book says so, one page before the Part VI you read:

\begin{bookquote}{book p.~379 / pdf p.~398}
\textit{Mathematics is creative and open-ended. By virtue of the use of conceptual metaphors and
conceptual blends, present mathematics can be extended to create new forms by importing structure
from one branch to another and by fusing mathematical ideas from different branches.}
\end{bookquote}

\subsection{The defensible version of your worry}

\begin{critbox}[title={What survives the correction --- and it is not nothing}]
The right conclusion is not \emph{``there is a ceiling.''} It is:
\begin{center}\itshape\large
Unintuitive mathematics is slower to invent, harder to teach,\\and more error-prone.
\end{center}
That is a real and expensive cost. Non-Euclidean geometry took roughly two millennia, and it took
that long largely because the parallel postulate \emph{felt} necessary --- which is an embodiment
effect. But a tax on the rate of discovery is not a wall across the road.
\end{critbox}
